\documentclass{davian}

%% --- extra packages for this document (the class already loads the essentials) ---
\usepackage{tikz}
\usepackage{lipsum} % filler text for the demo; remove in real papers

\runningtitle{A DAVIAN Lab Preprint}

\title{Draft Less, Verify More:\\ A DAVIAN Lab Preprint Template}

\authors{%
  Kang Eun Jeon\affmark{1,*}\quad
  Second Author\affmark{1}\quad
  Third Author\affmark{1,2}\quad
  Jaegul Choo\affmark{1,\dag}%
}
\affiliations{%
  \affmark{1}DAVIAN Lab, KAIST AI \qquad
  \affmark{2}Collaborating Institution%
}
\contributions{%
  \affmark{*}Equal contribution.\quad
  \affmark{\dag}Corresponding author.%
}
\correspondence{\texttt{eekejeon@gmail.com}}

\begin{document}
\maketitle

\begin{abstract}
This document is both a usage example and a visual showcase of the DAVIAN Lab
LaTeX template. It demonstrates the title block with numbered affiliations and
contribution markers, the section heading style, math, figures, tables, and
citations. Replace this abstract with your own and start writing. The template
targets single-column arXiv-style preprints and technical reports.
\end{abstract}

\section{Introduction}
\label{sec:intro}
Welcome to the DAVIAN Lab preprint template. The body text is set in XCharter,
section headings in Montserrat, and accents use the KAIST deep blue. You can
cite work in the usual way~\cite{vaswani2017attention,he2016deep} and cross
reference sections with \texttt{cleveref}, e.g.\ \Cref{sec:method}.

\lipsum[1]

\section{Method}
\label{sec:method}
Inline math such as $\mathbf{y} = \mathbf{W}\mathbf{x} + \mathbf{b}$ uses
Charter-keyed glyphs so it blends with the body. Display math:
\begin{equation}
  \label{eq:softmax}
  \mathrm{softmax}(\mathbf{z})_i = \frac{e^{z_i}}{\sum_{j=1}^{K} e^{z_j}},
  \qquad i = 1,\dots,K.
\end{equation}

\subsection{A subsection}
\lipsum[2]

\paragraph{A run-in paragraph heading.}
Paragraph headings are set in Montserrat bold and run into the text.

\section{Experiments}
\label{sec:exp}

\Cref{fig:placeholder} shows a placeholder figure and \Cref{tab:results}
a results table.

\begin{figure}[t]
  \centering
  \begin{tikzpicture}
    \draw[kaistblue, very thick, rounded corners]
      (0,0) rectangle (10,3.2);
    \node[align=center, text=kaistbluedark, font=\sffamily\large]
      at (5,1.6) {Figure placeholder\\[2pt]
      \footnotesize replace with \texttt{\textbackslash includegraphics}};
  \end{tikzpicture}
  \caption{A placeholder figure drawn with Ti\emph{k}Z so the template compiles
  with no external image files. Swap in your own graphic.}
  \label{fig:placeholder}
\end{figure}

\begin{table}[t]
  \centering
  \caption{Example results table using \texttt{booktabs} rules.}
  \label{tab:results}
  \begin{tabular}{lcc}
    \toprule
    Method & Accuracy ($\uparrow$) & Latency (ms, $\downarrow$) \\
    \midrule
    Baseline        & 71.2 & 18.4 \\
    + Improvement A  & 74.8 & 17.9 \\
    \textbf{Ours}    & \textbf{78.5} & \textbf{15.1} \\
    \bottomrule
  \end{tabular}
\end{table}

\lipsum[3]

\section{Conclusion}
We provided a clean, lab-branded single-column template. Edit \texttt{main.tex}
and \texttt{references.bib} to start your own paper.

\bibliographystyle{plainnat}
\bibliography{references}

\end{document}
